\section{Introduction}\label{sec:intro}

Language models increasingly make affect-laden subjective judgments: they moderate emotionally charged posts, co-author stories, hold companionship-style conversations, and critique creative work. All of these tasks hinge on a model's internal sense of how \emph{dramatic} a piece of text is---and on whether that sense matches the reader being served. But ``dramatic'' is not one thing. A war memoir and a soap-opera confession can be equally intense while differing in kind: one earns its intensity, the other piles it on. If a model collapses drama onto a single intensity axis, it conflates earned drama with melodramatic excess, and it mis-models readers who disagree not about \emph{how much} drama a text has but about \emph{what kind} they are willing to credit. Knowing whether a model's affect representation has the right shape---and whether reader disagreement lives in that shape---is a prerequisite for steering, auditing, or personalizing any of these systems.

\para{The problem.} Two questions are tangled together in practice. First, is the \emph{construct} of drama one-dimensional or many? Whether melodrama is ``just more drama'' or a qualitatively distinct register determines whether a single learned ``drama'' direction can faithfully control affective tone. Second, where does \emph{reader} disagreement come from? When two annotators split on whether a passage is dramatic, are they reading along a different internal direction (a difference in kind), or sharing the same axis and merely placing their cutoff in a different spot (a difference in degree)? These have opposite engineering consequences---personalization by re-aiming a direction versus personalization by shifting a threshold---yet they are routinely conflated.

\para{The gap.} No prior work studies ``drama'' as a model-internal construct, nor the dramatic-versus-melodramatic contrast. The closely related single-direction-versus-subspace question is unsettled even for the best-studied concept, refusal: \citet{arditi2024refusal} report a single mediating direction, while \citet{wollschlager2025geometry} argue for a multi-dimensional concept cone. Separately, work on emotion shows that affect occupies a low-dimensional manifold in representation space \citep{reichman2025emotions}, and work on annotation shows that reader effects are real and large \citep{orlikowski2025beyond,zhou2025subjectivity}. But model-internal affect geometry has never been joined to per-reader behavioral variation: we do not know whether the dimensions a model uses to represent affect are the same dimensions along which readers actually disagree.

\para{Our approach.} We separate the geometry of the drama \emph{construct} from the geometry of \emph{reader} variation, and we state two precise hypotheses. \textbf{H1 (single-axis limit):} the dramatic-versus-melodramatic distinction and reader variation in it are thresholds and regions on a low-dimensional shared drama-feature space---in the limit, one drama axis plus a per-reader threshold. \textbf{H2 (reader directions):} reader variation is a reader-specific \emph{direction} within a multi-dimensional drama subspace, so different readers read along different axes. We test these with difference-in-means directions \citep{marks2023geometry} over hidden states of \qwenbig (replicated in \qwensmall), using a controlled, blind-validated drama corpus (\dramacorpus) for construct geometry and a 70-reader \goemotions setup for behavioral disagreement. As a teaser, \figref{fig:excess} shows the key construct result: dramatic and melodramatic texts sit at the \emph{same} value on the drama axis but split along an orthogonal ``excess'' axis.

\para{Findings.} Our result splits cleanly across the two questions. The \emph{construct} is genuinely multi-dimensional, refuting the single-axis limit: drama itself decodes perfectly (AUROC $1.00$), yet melodrama is not just more drama---the dramatic$\rightarrow$melodramatic step rotates to $\sim\!83^\circ$ (nearly orthogonal) to the drama axis in deep layers, and an 8-subconcept subspace has participation ratio rising to $\sim\!5.6$ (versus $\sim\!8.0$ for random directions and $1$ for a single axis). \emph{Reader} variation, in sharp contrast, is overwhelmingly a threshold on a shared low-dimensional axis (H1). Adding a per-reader threshold to a shared single axis lifts held-out AUROC from $0.51$ to $0.62$ on the transferred corpus subspace and from $0.66$ to $0.70$ in domain---while letting each reader choose their \emph{own} direction adds exactly the shuffled-reader null ($\Delta$AUROC $+0.0076$ real versus $+0.0076$ null, $z\!\approx\!1.0$, n.s.). Six prompted reader personas leave the drama axis essentially fixed (pairwise cosine $\geq 0.97$, mean $0.99$). In one sentence: {\bf drama is multi-dimensional and shared across readers, but people mostly differ in \emph{where} they put the threshold, not in \emph{which direction} they read.}

\para{Contributions.} We make the following contributions:
\begin{itemize}[leftmargin=*,itemsep=0pt,topsep=0pt]
\item We build \dramacorpus, a validated controlled drama corpus (24 scenarios $\times$ 4 registers plus 8 drama subconcepts, blind re-labeled at $100\%$ agreement), and a per-reader \goemotions setup (70 raters, 56{,}995 texts, $36\%$ contested), jointly enabling construct and reader geometry to be measured on the same model internals.
\item We show the drama construct is multi-dimensional: drama decodes at AUROC $1.00$, but melodrama lives on an orthogonal excess axis ($\sim\!83^\circ$, subspace PR $\sim\!5.6$), so melodrama is not just more drama (\figref{fig:excess}).
\item We show reader variation is a \emph{threshold, not a direction}, via nested held-out logistic models against a shuffled-reader permutation null: the decisive lift is adding a per-reader threshold ($+0.04$ to $+0.11$), while reader-specific directions add nothing beyond the null, replicated across scales ($1.5$B/$7$B) and layers.
\item We corroborate with a persona probe: re-deriving the drama axis under six prompted reader personas leaves its direction unchanged (pairwise cosine $\geq 0.97$), shifting projection and threshold rather than axis.
\end{itemize}

We review related work in \secref{sec:related}, detail our corpus, models, and nested-model protocol in \secref{sec:method}, present the construct and reader-variation experiments in \secref{sec:results}, and discuss implications for steering and personalization in \secref{sec:discussion}.
